\chapter{Εισαγωγή}
\label{ch:introduction}

\section{Σημασία της υγρασίας εδάφους για το οικοσύστημα}


\par Ως υγρασία εδάφους, ορίζεται η ποσότητα του νερού που συγκρατείται μεταξύ των σωματιδίων του εδάφους κυρίως στο ενεργό στρώμα του, το οποίο αφορά το ανώτερο επίπεδο έως και 1-2m. Πρόκειται για μία μεταβλητή ιδιότητα του εδάφους στο χρόνο, που αλλάζει δυναμικά ανάλογα τις βροχοπτώσεις, την άρδευση, τα υπόγεια ύδατα, ενώ παράληλα εξαρτάται και από διεργασίες που καταναλώνουν νερό όπως η εξάτμιση, η διαπνοή και η απορροή. 

\par Η εδαφική υγρασία αποτελεί ένα από τα σημαντικότερα στοιχεία ενός οικοσυστήματος, καθώς αποτελεί απαραίτητη προυπόθεση για την ανάπτυξη κι επιβίωση  φυτών και οργανισμών. Συμβάλλει σε αναγκαίους για τη ζωή μηχανισμούς της φύσης, όπως η συσσώρευση ή απομάκρυνση οργανικών και ανόργανων ενώσεων, η φωτοσύνθεση, η μεταφορά θρεπτικών ουσιών, καθώς και η ρύθμιση της θερμοκρασίας των φυτών μέσω της διαπνοής. Επομένως, η επαρκής διαθεσιμότητα νερού στο έδαφος είναι ύψιστης σημασίας για την ομαλή ανάπτυξη των φυτών.

\par  Η έλλειψη ή η υπερβολική ποσότητα νερού προκαλεί διαταραχές στη φυσιολογική λειτουργία των φυτών. Η ανεπαρκής ποσότητα νερού στο έδαφος οδηγεί στον κλείσιμο των οπών των φυλλωμάτων για τη μέγιστη διατήρη του νερού, περιορίζοντας ωστόσο τη διάχυση CO2 διαμέσου αυτών και κατά συνέπεια τη φωτοσύνθεση. Αντίθετα, σε συνθήκες υπερβολικής παρουσίας νερού μειώνεται ο χώρος που είναι διαθέσιμος για αέρα στο έδαφος, με αποτέλεσμα να περιορίζεται η διάχυση του οξυγόνου προς τις ρίζες. 

\section{Καταστάσεις υγρασίας του εδάφους και σημασία τους για τα φυτά}

Ανάλογα με την ποσότητα νερού που υπάρχει στο έδαφος, διαμορφώνονται διαφορετικές συνθήκες που επηρεάζουν την πρόσληψη νερού από τα φυτά και τον αερισμό του εδάφους. Δύο σημαντικές καταστάσεις που σχετίζονται με την περιεκτικότητα του εδάφους σε νερό είναι η υδατοϊκανότητα και ο κορεσμός.

\subsection{Υδατοϊκανότητα}
Η υδατοϊκανότητα (Field Capacity) αναφέρεται στην κατάσταση κατά την οποία το ελεύθερο 
νερός έχει αποστραγγιστεί  από τους πόρους του έδαφος λόγω της βαρύτητας. Το έδαφος
στην περίπτωση αυτή περιέχει ικανοποιητική ποσότητα νερού ενώ διατηρεί και αέρα στους πόρους. Αποτελεί μία ευνοϊκή συνθήκη για την ανάπτυξη των φυτών καθώς συμβάλλει στην κάλυψη των αναγκών τους σε νερό χωρίς να περιορίζεται σημαντικά ο αερισμός της ριζόσφαιρας.

\subsection{Κορεσμός}
Αντίθετα, ο κορεσμός του εδάφους χαρακτηρίζεται από την πλήρωση όλων των πόρων από νερό. Η παρατεταμένη περίοδος του εδάφους σε κορεσμό, οδηγεί σε μειωμένη πρόσληψη οξυγόνου από τα φυτά, καθώς και την απομάκρυνση θρεπτικών συστατικών από τις ρίζες των φυτών. Επομένως, αποτελεί μια
κατάσταση που θα πρέπει να αποφεύγεται για την ευημερή ανάπτυξη καλλιεργιών. 

\subsection{Σημείο Μόνιμης Μάρανσης}
Το Σημείο Μόνιμης Μάρανσης είναι η συνθήκη κατά την οποία το νερό συγκρατείται από το σωματίδια του εδάφους με ισχυρούς δεσμούς, καθιστώντας σχεδόν αδύνατη την απορρόφησή του από τις ρίζες των φυτών. Πρόκειται για μια ιδιότητα του εδάφους που εξαρτάται από τον τύπο, ενώ διαφέρει και στα διάφορα φυτά ανάλογα την ανθεκτικότητα και τη δυνατότητά τους να απορροφήσουν νερό.

